% Einheit 2 von 3 – Zentrische Streckung und Ähnlichkeit (bank/strahlensaetze/e2.jsonl, zusammenbau v0.1)
%% TODO Zweigzeile Teil 1: was hier gelernt wird (Satz in Schülersprache aus dem Einheitstitel) – Einheit 2
\einheitenkopf[e2]{Einheit 2 von 3 · Zentrische Streckung und Ähnlichkeit}
\zweigzeile{ab Kl. 8, je nach Buch bis Kl. 9 (am Gymnasium ab Kl. 7) · keine P10-Aufgabe}
\setcounter{aufgabe}{15}

%% TODO Titel: Ich-kann-Satz fehlt in der Bank (Platzhalter: Kettenname) – Nr. 16
%% TODO Anweisung: ein Satz über der ersten Teilaufgabe fehlt in der Bank – Nr. 16
\begin{aufgabe}{Zentrische Streckung und Ähnlichkeit}
%% TODO \rechenplatz{n}: Schrittzahl des Grundfalls fehlt in der Bank (nur bei fester Schreibform, 2.3 b)
\begin{teile}
\teil Rechteck $2\,\text{cm} \times 3\,\text{cm}$ und Rechteck $4\,\text{cm} \times 6\,\text{cm}$ – gleiche Form? \kreuz{gleiche Form} \\ \kreuz{andere Form}
\teil Strecke das Dreieck $ABC$ vom Zentrum $Z = A$ aus mit dem Streckfaktor $k = 2$. Zeichne das Bilddreieck und gib die Bildpunkte an.

\begin{ksys}[xmin=0,xmax=10,ymin=0,ymax=7]\punkt{0}{0}{A}\punkt{3}{0}{B}\punkt{0}{2}{C}\end{ksys}
\teil Das Zentrum $Z$ liegt außerhalb des Dreiecks $ABC$. Zeichne Strahlen von $Z$ durch die Ecken und strecke mit $k = 2$. Gib die Bildpunkte an.

\begin{ksys}[xmin=0,xmax=10,ymin=0,ymax=7]\punkt{0}{0}{Z}\punkt{1}{1}{A}\punkt{3}{1}{B}\punkt{2}{3}{C}\end{ksys}
\teil Verkleinere das Dreieck $ABC$ vom Zentrum $Z$ aus mit $k = 0{,}25$. Zeichne das Bild und gib die Bildpunkte an.

\begin{ksys}[xmin=0,xmax=10,ymin=0,ymax=9]\punkt{0}{0}{Z}\punkt{4}{4}{A}\punkt{8}{4}{B}\punkt{8}{8}{C}\end{ksys}
\teil $\overline{AB} = 5\,\text{cm}$, die Bildstrecke $\overline{A'B'} = 15\,\text{cm}$ – wie groß ist der Streckfaktor $k$? k = \leerfeld
\teil Ein Dreieck wird mit $k = 3$ zentrisch gestreckt. Welche Aussage stimmt? \kreuz{Die Bildseiten sind parallel zu den Originalseiten.} \\ \kreuz{Die Winkel werden dreimal so groß.} \\ \kreuz{Jede Seite wird um $3\,\text{cm}$ länger.}
\teil Dreieck 1 hat die Winkel $50^\circ$ und $70^\circ$, Dreieck 2 die Winkel $70^\circ$ und $60^\circ$. Sind die Dreiecke ähnlich? \kreuz{ja} \\ \kreuz{nein}
\teil Dreieck $ABC$ hat die Seiten $4\,\text{cm}$, $6\,\text{cm}$ und $7\,\text{cm}$, Dreieck $DEF$ die Seiten $10{,}5\,\text{cm}$, $6\,\text{cm}$ und $9\,\text{cm}$. Paare die Seiten der Größe nach und prüfe die Verhältnisse. Sind die Dreiecke ähnlich?
\teil Die Dreiecke $ABC$ und $DEF$ sind ähnlich, $AB$ entspricht $DE$. Es ist $AB = 4\,\text{cm}$, $DE = 6\,\text{cm}$ und $BC = 7\,\text{cm}$. Berechne $EF$. EF = \leerfeld[cm]
\teil Die beiden rechtwinkligen Dreiecke sind ähnlich. Miss in beiden die Kathete $a$ und die Hypotenuse $c$ und berechne jeweils $a : c$. Was fällt dir auf?

\dreieckrw{2.4}{1.8}{a}{b}{c} \quad \dreieckrw{4.8}{3.6}{a}{b}{c}
\teil Ein Rechteck hat den Flächeninhalt $6\,\text{cm}^2$. Es wird mit $k = 3$ zentrisch gestreckt. Wie groß ist der Flächeninhalt des Bilds? \leerfeld[cm²]
\teil Ein Würfelmodell hat $2\,\text{cm}$ Kantenlänge und wird mit $k = 3$ vergrößert. Wie lang ist eine Kante, wie groß die Oberfläche des großen Würfels?
\teil Ein Foto ist $9\,\text{cm}$ breit und $13\,\text{cm}$ hoch. Poster A ist $27\,\text{cm}$ breit und $39\,\text{cm}$ hoch, Poster B $30\,\text{cm}$ breit und $40\,\text{cm}$ hoch. Welches Poster ist zum Foto ähnlich? Gib den Streckfaktor an. Ein Baum ist auf dem Foto $2{,}4\,\text{cm}$ hoch – wie hoch ist er auf diesem Poster?
\end{teile}
\end{aufgabe}

%% TODO Titel: Ich-kann-Satz fehlt in der Bank (Platzhalter: Kettenname) – Nr. 17
%% TODO Anweisung: ein Satz über der ersten Teilaufgabe fehlt in der Bank – Nr. 17
\begin{aufgabe}{Zentrum aus Original und Bild finden}
\begin{teile}
\teil Das Dreieck $A'B'C'$ ist das Bild von $ABC$ bei einer zentrischen Streckung. Zeichne Geraden durch $A$ und $A'$, $B$ und $B'$, $C$ und $C'$. Gib das Zentrum $Z$ an.

\begin{ksys}[xmin=0,xmax=9,ymin=0,ymax=7]\punkt{2}{2}{A}\punkt{3}{2}{B}\punkt{2}{3}{C}\punkt{3}{3}{A'}\punkt{5}{3}{B'}\punkt{3}{5}{C'}\end{ksys}
\end{teile}
\end{aufgabe}

%% TODO Titel: Ich-kann-Satz fehlt in der Bank (Platzhalter: Kettenname) – Nr. 18
%% TODO Anweisung: ein Satz über der ersten Teilaufgabe fehlt in der Bank – Nr. 18
\begin{aufgabe}{Zentrische Streckung und Ähnlichkeit – Fehler finden, Begründen, Anwendung}
\begin{teile}
\teil Die Originalstrecke ist $2\,\text{cm}$ lang, die Bildstrecke $7\,\text{cm}$. Tom rechnet: $k = 7 - 2 = 5$. Finde den Fehler und rechne richtig.
\teil Warum sind alle Quadrate zueinander ähnlich? Begründe.
\teil Ein Schullogo ist $6\,\text{cm}$ breit und $4\,\text{cm}$ hoch. Es soll ähnlich auf ein Banner vergrößert werden und dort $1{,}8\,\text{m}$ breit sein. Das Banner ist $1{,}5\,\text{m}$ hoch. Passt das Logo?
\end{teile}
\end{aufgabe}
